\chapter{Kategorik Veriler, Çapraz Tablolar ve Ki-Kare Testi}
\index{kategorik veri}
\index{çapraz tablo}
\index{ki-kare testi}

\begin{hedefler}
Çapraz tabloyu frekans ve yüzdelerle okumak; beklenen frekansları hesaplamak;
ki-kare testi ile Cramér $V$ sonucunu birlikte yorumlamak.
\end{hedefler}

\begin{hazirlik}
Grup büyüklükleri çok farklıysa yalnız hücre sayılarını karşılaştırmak neden
yanıltıcı olabilir? Satır veya sütun yüzdelerinin rolünü açıklayın.
\end{hazirlik}

\section{Kategorik veriyi özetlemek}
\index{frekans}
\index{satır yüzdesi}
\index{sütun yüzdesi}

Kategorik değişkenlerde temel özet frekans ve yüzdedir. İki kategorik
değişken birlikte incelendiğinde satır ve sütun kategorilerinden oluşan çapraz
tablo kullanılır. Grup büyüklükleri farklıysa yalnız hücre sayılarına değil,
uygun satır veya sütun yüzdelerine bakılmalıdır.

\section{Bağımsızlık hipotezi}
\index{bağımsızlık!kategorik değişkenler}
\index{beklenen frekans}

Ki-kare bağımsızlık testinde
\[
H_0:\text{iki kategorik değişken bağımsızdır}
\]
hipotezi sınanır. Bağımsızlık altında $(i,j)$ hücresinin beklenen frekansı
\[
E_{ij}=\frac{R_iC_j}{N}
\]
ile hesaplanır. Burada $R_i$ satır, $C_j$ sütun toplamıdır.

\section{Pearson ki-kare istatistiği}
\index{Pearson ki-kare testi}
\index{standartlaştırılmış artık}
\index{hücre katkısı}

Gözlenen ve beklenen frekanslar arasındaki standartlaştırılmış farklar
\[
\chi^2=\sum_i\sum_j\frac{(O_{ij}-E_{ij})^2}{E_{ij}}
\]
istatistiğinde birleştirilir. $r\times c$ tabloda serbestlik derecesi
$(r-1)(c-1)$'dir \cite{degroot2012,conover1999}.

Büyük $\chi^2$ değeri gözlenen tablonun bağımsızlık altında beklenen tablodan
uzak olduğunu gösterir. Hangi hücrelerin farkı oluşturduğunu anlamak için
standartlaştırılmış artıklar ve tablo yüzdeleri incelenir.

\section{Uygulama koşulları}
\index{ki-kare testi!beklenen frekans koşulu}
\index{Fisher kesin testi}

Gözlemler bağımsız olmalı ve beklenen frekanslar ki-kare yaklaşımı için
yeterince büyük olmalıdır. Aynı kişinin birden çok kez sayılması bağımsızlığı
bozar. Küçük $2\times2$ tablolarda Fisher kesin testi düşünülebilir.

\section{Etki büyüklüğü}
\index{Cramer V@Cramér $V$}
\index{etki büyüklüğü!kategorik veri}

İlişkinin büyüklüğü Cramér katsayısıyla
\[
V=\sqrt{\frac{\chi^2}{N\min(r-1,c-1)}}
\]
özetlenebilir. Anlamlı test ilişkinin güçlü veya uygulamada önemli olduğunu
tek başına göstermez.

\begin{outputguidebox}
Çıktıda Pearson ki-kare satırını, serbestlik derecesini ve $p$-değerini
okuyun. Beklenen frekans uyarısını, satır/sütun yüzdelerini, artıkları ve
Cramér $V$ değerini ayrıca inceleyin.
\end{outputguidebox}

\section{Raporlama örneği}

``Tercih edilen çözüm stratejisi ile öğretim yöntemi arasında ilişki bulundu,
$\chi^2(2,N=180)=8{,}31$, $p=0{,}016$, $V=0{,}21$. Yüzdeler ve
standartlaştırılmış artıklar farkın özellikle görsel strateji kategorisinde
yoğunlaştığını gösterdi.''

\section{Literatür veri setiyle IBM SPSS laboratuvarı}
\index{Student Performance veri seti!ki-kare testi}
\index{SPSS!Crosstabs komutu}
\index{okul!başarı çapraz tablosu}
\index{odds oranı!ki-kare uygulaması}
\index{Cramer V@Cramér $V$!SPSS uygulaması}

Bu uygulamada \emph{Student Performance} matematik dosyasındaki okul
değişkeni ile yıl sonu başarı durumu arasındaki ilişki incelenecektir
\cite{cortez2008data,cortezsilva2008}. Veri setinde okul iki kategoriyle
kodlanmıştır: Gabriel Pereira (GP) ve Mousinho da Silveira (MS). Yıl sonu
notu \texttt{G3} en az 10 olan öğrenciler başarılı, 10'un altında olanlar
başarısız olarak sınıflandırılacaktır.

Araştırma hipotezleri
\[
H_0:\text{okul ile başarı durumu bağımsızdır},
\]
\[
H_1:\text{okul ile başarı durumu bağımsız değildir}
\]
biçimindedir. Analiz birimi öğrencidir ve her öğrenci çapraz tabloda yalnız
bir hücreye katkıda bulunur.

\begin{researchbox}
Başarı durumunu iki kategoriye indirmek G3 puanları arasındaki büyüklük
bilgisini kaybettirir. Eşik eğitim sistemindeki anlamıyla önceden
gerekçelendirilmelidir. Ayrıca okul ile başarı arasında ilişki bulunması okulun
nedensel etkisini göstermez; öğrenci bileşimi, önceki başarı ve okul dışı
koşullar olası karıştırıcılardır.
\end{researchbox}

\subsection{SPSS syntax}

Başarı göstergesi önce sistem eksik olarak oluşturulur; böylece eksik G3
değerleri yanlışlıkla başarısız kategorisine atanmaz. \textbf{Crosstabs}
komutu gözlenen ve beklenen frekansları, satır ve sütun yüzdelerini, hücre
artıklarını, ki-kare testlerini, Cramér $V$ ve risk ölçülerini üretir.

\begin{lstlisting}[language={}]
DATASET ACTIVATE StudentMath.

COMPUTE basarili=$SYSMIS.
IF NOT MISSING(G3) basarili=(G3>=10).
VARIABLE LABELS basarili 'Yıl sonu notu en az 10'.
VALUE LABELS basarili 0 'Başarısız' 1 'Başarılı'.
VALUE LABELS school
 'GP' 'Gabriel Pereira'
 'MS' 'Mousinho da Silveira'.
VARIABLE LEVEL school basarili (NOMINAL).
EXECUTE.

CROSSTABS
 /TABLES=school BY basarili
 /FORMAT=AVALUE TABLES
 /STATISTICS=CHISQ PHI RISK
 /CELLS=COUNT EXPECTED ROW COLUMN
        RESID SRESID ASRESID
 /COUNT ROUND CELL.
\end{lstlisting}

Menüyle çalışmak için \textbf{Analyze $>$ Descriptive Statistics $>$
Crosstabs} yolunda okul satır, başarı durumu sütun değişkeni olarak seçilir.
\textbf{Statistics} bölümünde \textbf{Chi-square}, \textbf{Phi and Cramer's
V} ve gerekirse \textbf{Risk}; \textbf{Cells} bölümünde gözlenen, beklenen,
satır yüzdesi ve ayarlanmış standartlaştırılmış artıklar işaretlenir.

\subsection{Gözlenen çapraz tablo}

\begin{table}[htbp]
\centering
\caption{Okula göre yıl sonu başarı durumu.}
\begin{tabular}{@{}lrrrrr@{}}
\toprule
 & \multicolumn{2}{c}{Başarısız} & \multicolumn{2}{c}{Başarılı} & Toplam \\
\cmidrule(lr){2-3}\cmidrule(lr){4-5}
Okul & $n$ & Satır \% & $n$ & Satır \% & $n$ \\
\midrule
GP & 113 & 32,4 & 236 & 67,6 & 349 \\
MS & 17 & 37,0 & 29 & 63,0 & 46 \\
\midrule
Toplam & 130 & 32,9 & 265 & 67,1 & 395 \\
\bottomrule
\end{tabular}
\end{table}

Grupların büyüklükleri çok farklıdır: öğrencilerin 349'u GP, 46'sı MS
okulundadır. Bu nedenle yalnız başarılı öğrenci sayılarını, 236 ile 29'u,
karşılaştırmak yanıltıcıdır. Okul içindeki başarı yüzdeleri karşılaştırıldığında
GP için yüzde 67,6, MS için yüzde 63,0 elde edilir; gözlenen fark 4,6 yüzde
puandır.

\subsection{Beklenen frekanslar ve hücre artıkları}

Bağımsızlık altında beklenen frekanslar şöyledir:

\begin{table}[htbp]
\centering
\caption{Bağımsızlık altında beklenen frekanslar ve ayarlanmış artıklar.}
\begin{tabular}{@{}lrrrr@{}}
\toprule
 & \multicolumn{2}{c}{Başarısız} & \multicolumn{2}{c}{Başarılı} \\
\cmidrule(lr){2-3}\cmidrule(lr){4-5}
Okul & Beklenen & Ayarlı artık & Beklenen & Ayarlı artık \\
\midrule
GP & 114,861 & $-0{,}621$ & 234,139 & 0,621 \\
MS & 15,139 & 0,621 & 30,861 & $-0{,}621$ \\
\bottomrule
\end{tabular}
\end{table}

En küçük beklenen frekans 15,139'dur; dolayısıyla bu tabloda büyük örneklem
Pearson ki-kare yaklaşımının beklenen frekans koşulu sağlanmaktadır. Ayarlanmış
artıkların mutlak değerleri 2'nin oldukça altındadır. Hiçbir hücre bağımsızlık
altında beklenenden belirgin biçimde fazla veya az görünmemektedir.

\subsection{Ki-kare ve etki büyüklüğü çıktısı}

\begin{table}[htbp]
\centering
\caption{Okul ile başarı durumu için ilişki testleri ve etki büyüklüğü.}
\begin{tabular}{@{}lrrr@{}}
\toprule
Ölçü & Değer & $sd$ & $p$ \\
\midrule
Pearson ki-kare & 0,386 & 1 & 0,534 \\
Süreklilik düzeltmesi & 0,206 & 1 & 0,650 \\
Fisher kesin testi & --- & --- & 0,617 \\
Cramér $V$ & 0,031 & --- & 0,534 \\
\bottomrule
\end{tabular}
\end{table}

Pearson sonucu $\chi^2(1,N=395)=0{,}386$, $p=0{,}534$'tür. Yüzde 5
anlamlılık düzeyinde okul ile başarı durumu arasında ilişki bulunduğunu
söylemek için yeterli istatistiksel kanıt yoktur. Bu sonuç iki okulun başarı
oranlarının kesinlikle eşit olduğunu kanıtlamaz; gözlenen farkın örnekleme
değişkenliğiyle uyumlu olduğunu gösterir.

$2\times2$ tabloda Cramér $V$, mutlak phi katsayısına eşittir. $V=0{,}031$
örneklemde çok küçük bir ilişkiye işaret eder. Etki büyüklüğü ve güven aralığı,
yalnız $p$-değerine göre daha açık bir büyüklük değerlendirmesi sağlar.

\subsection{Odds oranı}

Başarı odds'u GP için $236/113$, MS için $29/17$'dir. GP'nin MS'ye göre
başarı odds oranı
\[
OR=\frac{236\times17}{113\times29}\approx1{,}224
\]
ve yaklaşık yüzde 95 güven aralığı $[0{,}646;2{,}320]$'dir. Aralık 1'i
içerdiğinden veriler daha düşük veya daha yüksek odds değerleriyle uyumludur;
tek bir yön ve büyüklük hakkında kesin bir sonuç vermez.

SPSS \textbf{Risk Estimate} tablosunda odds oranının yönü kategori sırasına
bağlıdır. Yazılımın hangi satır ve sütun kategorisini paya aldığını kontrol
etmeden yalnız sayıyı raporlamak ters yönlü yorum üretebilir. Bu nedenle
raporda ``GP'nin MS'ye göre başarı odds'u'' gibi yön açıkça yazılmalıdır.

\begin{mistakebox}
$p=0{,}534$ iki okulun aynı olduğunu kanıtlamaz. Özellikle MS grubunda yalnız
46 öğrenci bulunduğundan küçük veya orta büyüklükteki gerçek farklar hassas
biçimde tahmin edilemeyebilir. Eşitlik savı için yalnız anlamlı olmayan test
değil, önceden tanımlanmış eşdeğerlik sınırları ve yeterli örneklem gerekir.
\end{mistakebox}

\begin{outputguidebox}
Çıktıyı şu sırayla okuyun: (1) kategori kodları ve eksikler, (2) gözlenen
frekanslar, (3) araştırma sorusuna uygun satır veya sütun yüzdeleri,
(4) beklenen frekansların yeterliliği, (5) Pearson ki-kare ve gerekirse Fisher
testi, (6) ayarlanmış artıklarla farkın hangi hücrelerde yoğunlaştığı,
(7) Cramér $V$, (8) yönü açık odds veya risk oranı ve güven aralığı,
(9) gözlemsel tasarımın nedensellik sınırı. Yalnız ``Sig.'' değerini
raporlamayın.
\end{outputguidebox}

\subsection{Örnek bulgu paragrafı}

``GP okulundaki öğrencilerin yüzde 67,6'sı (236/349), MS okulundakilerin
yüzde 63,0'ı (29/46) yıl sonu notunda başarı eşiğini karşılamıştır. Okul ile
başarı durumu arasında istatistiksel olarak belirgin bir ilişki bulunmamıştır,
$\chi^2(1,N=395)=0{,}386$, $p=0{,}534$, Cramér $V=0{,}031$. Bütün beklenen
frekanslar 15'in üzerindedir ve ayarlanmış artıkların mutlak değeri 0,621'i
aşmamıştır. GP'nin MS'ye göre başarı odds oranı 1,22'dir (yüzde 95 GA
$[0{,}65;2{,}32]$). Sonuç eşitlik kanıtı veya okulun nedensel etkisi olarak
yorumlanmamıştır.''

\section{Bölüm sonu çözümlü problemler}

\begin{enumerate}
  \item \textbf{Beklenen frekans.} 120 öğrencilik bir tabloda program satır toplamı 50, ``başvurdu'' sütun toplamı 36'dır. Bağımsızlık altında bu hücrenin beklenen frekansını bulun.

  \textbf{Çözüm:} $E=(50\cdot36)/120=15$'tir. Gözlenen sayı bu değerden ayrıldıkça hücrenin ki-kare istatistiğine katkısı artar.

  \item \textbf{Ki-kare katkısı.} Önceki hücrede gözlenen frekans 21 ise hücre katkısını hesaplayın.

  \textbf{Çözüm:} Katkı $(O-E)^2/E=(21-15)^2/15=36/15=2{,}40$'tır. Toplam $\chi^2$, bütün hücre katkılarının toplamıdır.

  \item \textbf{Cramér V.} $2\times3$ tabloda $\chi^2=12$, $N=150$ olsun. Cramér $V$ değerini bulun.

  \textbf{Çözüm:} $\min(r-1,c-1)=\min(1,2)=1$ olduğundan $V=\sqrt{12/(150\cdot1)}=\sqrt{0{,}08}\approx0{,}283$'tür. Uygulamalı önem bağlama göre değerlendirilir.
\end{enumerate}

\bolumkaynaklari{\cite{degroot2012,conover1999,wasserman2004,cortez2008data,cortezsilva2008}}

\section*{Hatırla--Uygula--Yansıt}

\begin{enumerate}
  \item Bir çapraz tabloda neden yüzdelerin de incelenmesi gerektiğini açıklayın.
  \item Beklenen frekans formülünü kullanarak bir hücre için hesaplama yapın.
  \item Anlamlı genel testten sonra artıkların rolünü belirtin.
\end{enumerate}